\section{Nonlinear compatibility checks with explicit scope}
\label{sec:networkcompatibility}

The value equality in Proposition~\ref{prop:networklinear} makes a
behavior-valued task essential for testing source independence. A
compatibility checker must report either a constructive realization, a
violated necessary condition, or an unresolved case. Satisfying a finite
list of necessary conditions does not establish compatibility.

\subsection{Independence of source-disjoint marginals}
For a set $U$ of parties, let $E(U)$ denote the union of sources incident to
at least one party in $U$.
\begin{proposition}[Source-disjoint groups]
\label{prop:sourcedisjoint}
Suppose disjoint party groups $U,V$ satisfy $E(U)\cap E(V)=\varnothing$.
In an independent-source classical or tensor-product quantum network,
with no additional global shared randomness or communication,
\[
 p(a_U,a_V\mid x)=p(a_U\mid x_U)\,p(a_V\mid x_V).
\]
The equation is a marginal over all other outputs, not a condition on their
values. Inputs are fixed independently of the sources.
\end{proposition}
\begin{proof}
In the classical model, summing other local outputs removes their normalized
response factors. The remaining response factors depend on disjoint
collections of independent source variables, so their expectation factors.
In the quantum model, summing the other POVMs gives identities. Tracing their
systems out of the product source state leaves a product between the systems
in $U$ and those in $V$, since no source meets both groups. The Born trace
therefore factors. Independence from inputs outside each group follows in
the same way.
\end{proof}
In a bilocal chain this forces $p(a,c\mid x,z)=p(a\mid x)p(c\mid z)$.
Conditioning on Bob's output can correlate the end parties through
entanglement swapping, so the unconditional factorization must not be
applied to such postselected data. An observed violation excludes both the
classical and the specified quantum independent-source model, under the
stated assumptions.

\subsection{A triangle event inequality}
For the triangle with independent pair sources $AB$, $BC$, and $CA$, every
classical behavior obeys, at each fixed input tuple,
\begin{equation}
 p(a,b,c\mid x,y,z)^2\leq
 p_A(a\mid x)\,p_B(b\mid y)\,p_C(c\mid z).
 \label{eq:trianglefinner}
\end{equation}
This is the triangle specialization of Finner's product-space inequality
\cite{finner1992}. Network versions and a quantum extension for independent
bipartite sources are established by Renou et al.~\cite{renou2019limits}.
The checker records applicability to both classical and quantum independent
pair sources. The elementary proof below establishes the classical case;
the quantum case uses the cited theorem rather than replacing operator
algebra by classical hidden variables.

Let $r,s,t$ be the independent variables of the three sources and let
$f(r,t)$, $g(r,s)$, and $h(s,t)$ be the local conditional probabilities of
the selected events. Private local randomness is already absorbed into
these functions, so $0\leq f,g,h\leq1$. At the fixed inputs,
$p(a,b,c)=\mathbb E[fgh]$. Two Cauchy--Schwarz applications give
\begin{align*}
 \mathbb E[fgh]
 &\leq \mathbb E_{s,t}\!\left[h(s,t)
       \sqrt{\mathbb E_r f(r,t)^2\,\mathbb E_r g(r,s)^2}\right],\\
 \bigl(\mathbb E[fgh]\bigr)^2
 &\leq \mathbb E[h^2]\,
   \mathbb E_{s,t}\!\left[\mathbb E_r f(r,t)^2\,
                         \mathbb E_r g(r,s)^2\right]\\
 &=\mathbb E[f^2]\,\mathbb E[g^2]\,\mathbb E[h^2]
 \leq\mathbb E[f]\,\mathbb E[g]\,\mathbb E[h].
\end{align*}
The last expectations are the three marginal event probabilities.

For example, the common-random-bit distribution
$p(000)=p(111)=1/2$ violates Eq.~\eqref{eq:trianglefinner}:
the event $000$ has squared probability $1/4$, while the product of its
marginals is $1/8$. Thus a common tripartite classical source cannot be
silently replaced by three independent pair sources. A nonviolation of
these event inequalities supplies no general realization theorem. The
checker does not interpret an arbitrary nonsignalling table as a network
of independent nonsignalling resources; that is a separate physical model.

\subsection{An exactly solvable no-input chain fragment}
When all parties in a bilocal chain have one input, source-disjoint
independence is also sufficient for compatibility, for arbitrary finite
output alphabets \cite{renou2019limits}.
\begin{proposition}[No-input bilocal reconstruction]
\label{prop:noinputbilocal}
A normalized distribution $p(a,b,c)$ admits a classical independent-source
realization on $A$--$B$--$C$ if and only if
$p_{AC}(a,c)=p_A(a)p_C(c)$. The same condition characterizes the
finite-dimensional quantum independent-source behaviors in this no-input
fragment.
\end{proposition}
\begin{proof}
Necessity in both models follows from
Proposition~\ref{prop:sourcedisjoint}. For sufficiency, the first source
samples $\lambda=a$ from $p_A$, and the second samples $\nu=c$ from $p_C$,
independently. Alice and Charlie output these variables. Bob receives both
and samples
\[
 q(b\mid\lambda=a,\nu=c)=
 \frac{p(a,b,c)}{p_A(a)p_C(c)}
\]
whenever the denominator is positive. The assumed factorization normalizes
this response. If the denominator is zero, normalization and nonnegativity
force every $p(a,b,c)$ in that row to vanish, and Bob may use any normalized
response there. The constructed joint probability is exactly $p(a,b,c)$.
A finite classical source/response model embeds into quantum theory using
orthogonal classical registers and diagonal POVMs.
\end{proof}
Rational input data produce rational source probabilities and responses, so
both the independence condition and the reconstructed witness can be checked
exactly. The proposition does not extend to multi-input bilocal games: one
must then use a common source model across the input choices, and the
marginal equality alone is insufficient.

\subsection{Certificates, tolerances, and unresolved cases}
For rational tables, a violation certificate can name an input profile, an
output event or marginal event, and the exact rational residual. A checker
recomputes that residual from the full supplied table and source graph,
including normalization and all applicability conditions. No floating square
root is needed for Eq.~\eqref{eq:trianglefinner}. Such a certificate excludes
a resource model for the supplied distribution; it is not an experimental
confidence statement.

The general compatibility checks accept only exact rational, normalized,
non-signalling tables and reject floating input. The separate bilocal
inequality evaluator also supports floating tables with explicit tolerances
and numerical evidence labels. No violation found means only that the
implemented conditions did not exclude the data. The constructive no-input chain result is an exception with a
proved sufficient condition and an explicit witness, not a precedent for
promoting generic passing checks to membership certificates.
